\section{Finite-sample likelihood geometry for adaptive experiments}
\label{sec:finite-sample-geometry}
This section supplies estimates at shrinking net resolution, not just a
fixed-resolution strong law. It also distinguishes a uniform probability
bound from a single event asserted for uncountably many policies.

Let \(\mathcal K\) be a compact metric parameter space with metric \(d\).
At stage \(i\), let \(\mathcal F_{i-1}\) be the observed past and
\(\mathcal A_i\) the information after all current actions but before the
readout. The noise \(\xi_i\), conditionally on \(\mathcal A_i\), is
\(N(0,\sigma^2)\), independent of that information. Policy kernels are
parameter independent on every history; candidate likelihoods hold the
entire realized action record fixed. Let \(f_i(\beta_0)=0\) and assume,
pathwise on all histories,
\begin{equation}
 |f_i(\beta)|\le M,\qquad
 |f_i(\beta)-f_i(\gamma)|\le d(\beta,\gamma),\qquad M\ge1.
 \label{eq:uniform-field-hypotheses}
\end{equation}
Each field is \(\mathcal A_i\)-measurable. Denote the logarithmic covering
number by \(\mathcal H_d(\epsilon)=\log N(\mathcal K,d,\epsilon)\).
Define
\[
 Z_n=n^{-1}\sum_i|\xi_i|,\quad z=3\sigma+1,\quad
 V=\max\{1,M^4,\sigma^2M^2\},\quad
 e(t)=\min\{1,t/(8M),t/(2z)\}.
\]

\begin{proposition}[Policy-uniform shrinking-net concentration]
\label{prop:finite-n-field-concentration}
For every \(n\ge1\), \(t>0\), and every admissible policy, including
horizon-dependent policies, the probability that either
\[
 \sup_\beta\left|n^{-1}\sum_i\xi_if_i(\beta)\right|>t
 \quad\text{or}\quad
 \sup_\beta\left|n^{-1}\sum_i\{f_i(\beta)^2-
 E[f_i(\beta)^2\mid\mathcal F_{i-1}]\}\right|>t
\]
is at most
\begin{equation}
 \Psi_n(t)=4\exp\{\mathcal H_d(e(t))-nt^2/(8V)\}+e^{-n}.
 \label{eq:finite-n-field-bound}
\end{equation}
All suprema are measurable. The same estimate holds uniformly over the
true parameter whenever the constants and covering numbers are common.
\end{proposition}
\begin{proof}
At one deterministic net point, the centered square is a martingale
difference bounded in absolute value by \(M^2\). Conditional exponential
bounds and Markov's inequality give
\[
 P\{|n^{-1}\sum_i(f_i^2-E_{i-1}f_i^2)|>t/2\}
 \le2e^{-nt^2/(8M^4)}.
\]
For the score, conditioning first on \(\mathcal A_i\) gives
\(E[e^{\lambda\xi_if_i}\mid\mathcal A_i]
 =e^{\lambda^2\sigma^2f_i^2/2}\).
Iterated conditioning therefore gives the corresponding bound
\(2e^{-nt^2/(8\sigma^2M^2)}\). These calculations do not condition on
future actions or on the final design matrix.

The elementary bound \(E e^{|\xi_i|/\sigma}\le2e^{1/2}\), followed by
iterated conditioning, proves
\(P(Z_n>3\sigma)\le e^{n(\log2+1/2-3)}\le e^{-n}\).
For parameters in one \(e(t)\)-ball, square-field interpolation costs
at most \(4Me(t)\le t/2\), including the conditional expectation;
score interpolation costs at most \(e(t)Z_n\le t/2\) on \(Z_n\le z\).
A union bound over the deterministic net proves \eqref{eq:finite-n-field-bound}.
Continuous fields on a compact metric space have suprema over a fixed
countable dense set, proving measurability. All constants are independent
of the policy and sample size.
\end{proof}

\subsection{An exact integrated-likelihood remainder}
Suppose the observation model is
\[
 Y_i=m_i^{\beta_0}+b_i^{\beta_0}(z_0)+\xi_i,
 \qquad f_i(\beta)=m_i^\beta-m_i^{\beta_0},
\]
and \(\sup_{\beta,\|z\|\le R}|b_i^\beta(z)|\le r_i\), for a common
deterministic summable sequence. Put
\(R_1=\sum_i r_i\), \(R_2=\sum_i r_i^2\).
The working prior \(\nu_n\) may be any probability on this Hilbert ball.
No density or weak convergence of \(\nu_n\) is assumed.

\begin{lemma}[Source-independent transient budget]
\label{lem:finite-n-exact-correction}
Writing \(\ell_n(\beta)=n^{-1}\log
[L_n^{\nu_n}(\beta)/L_n^{\nu_n}(\beta_0)]\), one has
\begin{equation}
 \ell_n(\beta)=\frac1{n\sigma^2}\sum_i\xi_if_i(\beta)
 -\frac1{2n\sigma^2}\sum_i f_i(\beta)^2+\mathcal R_n(\beta)/n,
 \qquad \sup_\beta|\mathcal R_n(\beta)|\le B_n,
 \label{eq:finite-n-exact-expansion}
\end{equation}
where
\[
 B_n=\sigma^{-2}\{2\sum_{i\le n}r_i|\xi_i|+3MR_1+3R_2\},\quad
 b_0=3(MR_1+R_2)/\sigma^2,\quad b_1=2R_1/\sigma.
\]
In particular, for \(u>0\),
\begin{equation}
 P\{B_n>b_0+b_1\sqrt{2u}\}\le2e^{-u}.
 \label{eq:finite-n-transient-tail}
\end{equation}
\end{lemma}
\begin{proof}
Let \(L_n^\star\) be the product Gaussian likelihood with means
\(m_i^\beta\). Gaussian completion expresses \(L_n^{\nu_n}/L_n^\star\)
as the integral of \(\exp V_n^\beta(z)\), where
\[
 V_n^\beta(z)=\sigma^{-2}\sum_i
 \{b_i^\beta(z)(Y_i-m_i^\beta)-\tfrac12 b_i^\beta(z)^2\}.
\]
Its absolute value is bounded by
\(\sigma^{-2}\{\sum r_i|\xi_i|+MR_1+3R_2/2\}\).
A normalized positive integral has logarithm between the infimum and
supremum of its exponent. Apply this at \(\beta\) and \(\beta_0\);
the ideal likelihood expansion has the additional term
\(\sigma^{-2}\sum b_i^{\beta_0}(z_0)f_i(\beta)\), bounded by
\(MR_1/\sigma^2\). This proves the stated budget.
For \(R_1>0\), convexity, with weights \(r_i/R_1\) and any remaining
weight at zero, gives
\[
 E\exp\{\lambda\sum_{i\le n}r_i|\xi_i|\}
 \le2\exp\{\lambda^2\sigma^2R_1^2/2\}.
\]
Optimizing Markov's inequality yields \eqref{eq:finite-n-transient-tail}.
For \(R_1=0\) the remainder is identically zero.
\end{proof}

\begin{theorem}[A finite-n posterior separation inequality]
\label{thm:finite-n-posterior-separation}
Let \(F\subset\mathcal K\) be a measurable set satisfying, pathwise,
\[
 \inf_{\beta\in F}n^{-1}\sum_iE[f_i(\beta)^2\mid\mathcal F_{i-1}]
 \ge q>0.
\]
If the prior mass of the \(d\)-ball of radius \(\sqrt q/4\) around the
truth is at least \(e^{-\mathcal P_d(\sqrt q/4)}\), then
\begin{equation}
 \Pi_n^{\nu_n}(F)\le
 \exp\{\mathcal P_d(\sqrt q/4)-nq/(4\sigma^2)\}
 \label{eq:finite-n-posterior-mass}
\end{equation}
except on an event of probability at most
\begin{equation}
 \Psi_n(q/16)+P\{B_n>nq/(32\sigma^2)\},
 \label{eq:finite-n-posterior-failure}
\end{equation}
where the second term is explicitly bounded by
\eqref{eq:finite-n-transient-tail}.
Thus \(\mathcal P_d\) is an \emph{upper} bound on minus log prior mass.
\end{theorem}
\begin{proof}
On the good event of Proposition~\ref{prop:finite-n-field-concentration}
with \(t=q/16\), the numerator likelihood ratio over \(F\) is at most
\(\exp\{-13nq/(32\sigma^2)+B_n\}\).
On the indicated prior ball, \(|f_i|\le\sqrt q/4\) at every stage, and
the same score bound gives a likelihood ratio at least
\(\exp\{-3nq/(32\sigma^2)-B_n\}\).
Divide the integrated numerator by the integrated denominator and use
\(2B_n\le nq/(16\sigma^2)\). The remaining exponent is
\(-nq/(4\sigma^2)+\mathcal P_d\), as claimed. Every coefficient in this
calculation is retained; there is no unspecified constant multiplying a
shrinking separation.
\end{proof}

\subsection{Triangular experiments and nonstationary Laplace limits}
For a fixed infinite policy, the usual noise coupling is available. For
horizon-dependent policies, choose any joint probability space carrying
one admissible row of length \(n\) for each \(n\); rows need not be
independent. For each fixed \(t>0\), the right side of
\eqref{eq:finite-n-field-bound} is summable in \(n\). So is the transient
probability at a fixed positive multiple of \(n\). Borel--Cantelli at
countably many \(t\downarrow0\) proves, for each such chosen policy
sequence and coupling,
\begin{equation}
 \|\ell_n+Q_n\|_\infty\longrightarrow0\quad\text{almost surely},\qquad
 Q_n(\beta)=\frac1{2n\sigma^2}\sum_i
 E[f_i(\beta)^2\mid\mathcal F_{i-1}].
 \label{eq:adaptive-uniform-likelihood-tracking}
\end{equation}
The statement is not an intersection of probability-one events over an
uncountable collection of policies. Uniformity means that the finite-n
bounds preceding it are common to the full class.

\begin{theorem}[Adaptive posterior Laplace tracking and cluster LDPs]
\label{thm:adaptive-laplace-tracking}
Assume the preceding hypotheses and a fixed full-support prior on
\(\mathcal K\). For every continuous \(\varphi:\mathcal K\to\R\),
\begin{equation}
 -\frac1n\log\int e^{-n\varphi(\beta)}\Pi_n^{\nu_n}(d\beta)
 -\inf_{\beta\in\mathcal K}\{\varphi(\beta)+Q_n(\beta)\}
 \longrightarrow0\quad\text{almost surely}.
 \label{eq:adaptive-laplace-tracking}
\end{equation}
The sequence \(Q_n\) is precompact in the uniform norm. Every uniform
cluster limit \(I\) is a good rate function, vanishes at \(\beta_0\),
and the corresponding subsequence of posteriors satisfies the full
open-set lower and closed-set upper LDP with rate \(I\) and speed \(n\).
If \(Q_n\to I\) uniformly along the full sequence, this is a full adaptive
posterior LDP. These statements hold on one event for all continuous
\(\varphi\) and all Borel sets.
\end{theorem}
\begin{proof}
One has \(0\le Q_n\le M^2/(2\sigma^2)\), \(Q_n(\beta_0)=0\), and
\(|Q_n(\beta)-Q_n(\gamma)|\le M d(\beta,\gamma)/\sigma^2\).
Arzela--Ascoli gives precompactness. For every radius \(s>0\), full
support and a finite \(s/2\)-cover give
\[
 m(s):=\inf_{\beta\in\mathcal K}\Pi\{d(\gamma,\beta)<s\}>0.
\]
For the equicontinuous family \(\varphi+Q_n\), integration over an
\(s\)-ball about a minimizer gives
\[
 \inf(\varphi+Q_n)\le -n^{-1}\log\int e^{-n(\varphi+Q_n)}d\Pi
 \le\inf(\varphi+Q_n)+\omega_\varphi(s)+Ms/\sigma^2
       -n^{-1}\log m(s).
\]
Use this direct upper/lower integral bound also
with \(\varphi=0\), use \eqref{eq:adaptive-uniform-likelihood-tracking},
then let \(n\to\infty\) and \(s\downarrow0\).
The argument is deterministic once the uniform tracking event holds, so it
holds there for every continuous function without another exceptional set.
On a subsequence with \(Q_n\to I\), it gives the uniform log-likelihood
limit \(-I\). For an open set integrate a small neighbourhood of any one
of its points; for a closed set take the supremum of the likelihood over
that compact set. The denominator has exponent zero by the preceding
argument. These give the two LDP bounds simultaneously for all Borel sets.
\end{proof}
